\section{8.3 중간고사 대비 총정리}
\begin{frame}{새 개념이 없다 --- ``재현 가능한가''를 점검한다}
  \begin{itemize}
    \item Block A(Week 2--8)의 마지막 블록. 새 개념 \textbf{없음}
      --- 지금까지의 ``알고 있는 것''이 \textbf{시험장에서 빈
      종이에 스스로 재현}될 수 있는지 점검.
    \item ``수업에선 이해했다''와 ``시험장에서 유도할 수 있다''는
      서로 다른 역량 --- 손유도와 ``왜?'' 질문에서 후자가 직접
      점수로 연결.
  \end{itemize}
  \vskip0.5em
  \begin{block}{이 절의 구조}
    시험에서 가장 자주 나오는 질문 유형 \textbf{4개}를 뽑고,
    그중 \textbf{``손으로 써야 하는'' 3가지(유도 템플릿 ①②③)}를
    실제로 종이에 풀고, 나머지는 자가진단 플로우와 확인 문제
    9개로 채운다. 실습 노트북은 세 템플릿의 숫자를 코드로
    검증.
  \end{block}
\end{frame}
